\documentclass[11pt]{article}
\usepackage[margin=1in]{geometry}
\usepackage{booktabs}
\usepackage{caption}
\usepackage{amsmath}
\usepackage{float}
\usepackage{graphicx}
\usepackage{adjustbox}

\begin{document}

\section*{Asymmetric Demand Multipliers by Shock Size: v2}

This note reports the v2 asymmetric specification. The regression is run
separately on positive-demand and negative-demand subsamples, dropping
opposite-sign observations rather than setting them to zero. The shock-size
bins use the pooled monthly standard deviation thresholds from
\texttt{reg\_table\_static.RDS}; that is, the same $\sigma$ thresholds are used
for positive and negative shocks within a demand type.

For each demand type and month, define
\[
d^{(1)}_{n,t} = d_{n,t}\mathbf{1}\{|d_{n,t}| < \sigma_t\}, \quad
d^{(2)}_{n,t} = d_{n,t}\mathbf{1}\{|d_{n,t}| \in [\sigma_t,2\sigma_t]\}, \quad
d^{(3)}_{n,t} = d_{n,t}\mathbf{1}\{|d_{n,t}| > 2\sigma_t\}.
\]
The Fama--MacBeth regression is
\[
r_{n,t} = \alpha_t + M_1 d^{(1)}_{n,t}
        + M_2 d^{(2)}_{n,t}
        + M_3 d^{(3)}_{n,t}
        + \Gamma X_{n,t} + \varepsilon_{n,t},
\]
estimated separately for BMI, FIT, and OFI demand. Columns (1)--(3) use BMI,
columns (4)--(6) use FIT, and columns (7)--(9) use OFI. Within each demand
type, the three columns correspond to no controls, predictor controls, and
predictor plus liquidity controls. BMI specifications also include the BMI
control set used in the main price-impact tables.

\paragraph{Takeaway.}
The positive-demand domain shows strong concavity: multipliers are larger for
small shocks and smaller for large shocks, especially for FIT and OFI. This is
not consistent with a short-sales-constraint interpretation under which the
multiplier should rise with shock size. On the negative-demand side, OFI
continues to show strong concavity. FIT provides weaker evidence of concavity:
the largest-bin multiplier is below the smallest-bin multiplier, although the
differences are not statistically precise. BMI estimates are noisy. Crucially,
the results do not show evidence that $M$ increases with shock size, so the
evidence is not consistent with leverage constraints generating larger
multipliers for larger shocks.

\begin{table}[H]
\centering
\caption{v2 shock-size bin summary statistics}
\label{tab:v2_bin_summary}
\begin{tabular}{llrrrr}
\toprule
Shock & Type & Bin & Obs & Mean $d_{n,t}$ & SD $d_{n,t}$ \\
\midrule
Negative & BMI & 1 & 3,922 & -0.001889 & 0.001698 \\
Negative & BMI & 2 & 472 & -0.011188 & 0.003155 \\
Negative & BMI & 3 & 287 & -0.020162 & 0.006973 \\
Negative & FIT & 1 & 176,791 & -0.001351 & 0.001312 \\
Negative & FIT & 2 & 40,711 & -0.005666 & 0.002131 \\
Negative & FIT & 3 & 17,203 & -0.011023 & 0.004662 \\
Negative & OFI & 1 & 247,374 & -0.013672 & 0.011270 \\
Negative & OFI & 2 & 40,829 & -0.056677 & 0.017312 \\
Negative & OFI & 3 & 18,668 & -0.137339 & 0.070365 \\
\addlinespace
Positive & BMI & 1 & 3,973 & 0.002133 & 0.001837 \\
Positive & BMI & 2 & 654 & 0.010492 & 0.003105 \\
Positive & BMI & 3 & 392 & 0.020672 & 0.006904 \\
Positive & FIT & 1 & 233,213 & 0.001444 & 0.001496 \\
Positive & FIT & 2 & 48,811 & 0.007076 & 0.002600 \\
Positive & FIT & 3 & 25,849 & 0.015667 & 0.008279 \\
Positive & OFI & 1 & 187,161 & 0.012776 & 0.011188 \\
Positive & OFI & 2 & 24,051 & 0.056741 & 0.017554 \\
Positive & OFI & 3 & 9,661 & 0.134889 & 0.054641 \\
\bottomrule
\end{tabular}
\end{table}

\begin{table}[H]
\centering
\caption{v2 asymmetric multipliers by shock-size bin}
\label{tab:v2_asym_multipliers}
\begin{adjustbox}{max width=\textwidth}
\begingroup
\renewcommand{\vspace}[1]{}
\begin{tabular}{lccccccccc}
  \hline \multicolumn{10}{c}{Panel A: Positive shocks} \\
 \hline
  & \multicolumn{3}{c}{BMI} & \multicolumn{3}{c}{FIT} & \multicolumn{3}{c}{OFI} \\ 
  \cmidrule(l){2-4} \cmidrule(l){5-7} \cmidrule(l){8-10} 
  & (1) & (2) & (3) & (4) & (5) & (6) & (7) & (8) & (9) \\ 
  $M_{\{|d_{n,t}| < \sigma\}}$ & $2.49^{*}$ & $2.53^{*}$ & $2.73^{*}$ & $8.74^{***}$ & $6.75^{***}$ & $6.43^{***}$ & $2.44^{***}$ & $2.42^{***}$ & $2.99^{***}$ \\ 
   & $(1.38)$ & $(1.48)$ & $(1.51)$ & $(1.67)$ & $(1.28)$ & $(1.03)$ & $(0.21)$ & $(0.16)$ & $(0.15)$ \\ 
  $M_{\{|d_{n,t}| \in [\sigma, 2\sigma]\}}$ & $0.88$ & $0.85$ & $1.05$ & $6.13^{***}$ & $5.05^{***}$ & $5.13^{***}$ & $1.95^{***}$ & $1.92^{***}$ & $2.27^{***}$ \\ 
   & $(1.08)$ & $(1.07)$ & $(1.13)$ & $(0.97)$ & $(0.76)$ & $(0.68)$ & $(0.18)$ & $(0.15)$ & $(0.13)$ \\ 
  $M_{\{|d_{n,t}| > 2 \sigma\}}$ & $1.16^{**}$ & $1.11^{**}$ & $1.09^{*}$ & $3.90^{***}$ & $3.44^{***}$ & $3.57^{***}$ & $1.45^{***}$ & $1.40^{***}$ & $1.56^{***}$ \\ 
   & $(0.58)$ & $(0.50)$ & $(0.58)$ & $(0.59)$ & $(0.53)$ & $(0.51)$ & $(0.17)$ & $(0.15)$ & $(0.14)$ \\ \vspace{5pt} 
  Predictor controls & N & Y & Y & N & Y & Y & N & Y & Y \\ 
  Liquidity controls & N & N & Y & N & N & Y & N & N & Y \\ \vspace{5pt} 
  Obs & 5,019 & 5,019 & 5,019 & 307,873 & 307,873 & 307,873 & 220,873 & 220,873 & 220,873 \\ 
  $R^2$ & $0.064$ & $0.184$ & $0.234$ & $0.012$ & $0.072$ & $0.090$ & $0.031$ & $0.101$ & $0.143$ \\ 
  Marginal $R^2(d_{n,t})$ & $0.024$ & $0.021$ & $0.022$ & $0.012$ & $0.008$ & $0.006$ & $0.031$ & $0.026$ & $0.029$ \\ 
 \hline
  $M_{\{|d_{n,t}| \in [\sigma, 2\sigma]\}} - M_{\{|d_{n,t}| < \sigma\}}$ & $-1.61$ & $-1.69$ & $-1.69$ & $-2.62^{***}$ & $-1.70^{**}$ & $-1.30^{**}$ & $-0.49^{***}$ & $-0.50^{***}$ & $-0.72^{***}$ \\ 
   & $(1.02)$ & $(1.13)$ & $(1.09)$ & $(0.93)$ & $(0.74)$ & $(0.65)$ & $(0.10)$ & $(0.10)$ & $(0.10)$ \\ 
  $M_{\{|d_{n,t}| > 2 \sigma\}} - M_{\{|d_{n,t}| \in [\sigma, 2\sigma]\}}$ & $0.28$ & $0.26$ & $0.05$ & $-2.22^{***}$ & $-1.61^{***}$ & $-1.56^{***}$ & $-0.49^{***}$ & $-0.52^{***}$ & $-0.71^{***}$ \\ 
   & $(0.74)$ & $(0.80)$ & $(0.81)$ & $(0.76)$ & $(0.60)$ & $(0.60)$ & $(0.12)$ & $(0.11)$ & $(0.11)$ \\ 
  $M_{\{|d_{n,t}| > 2 \sigma\}} - M_{\{|d_{n,t}| < \sigma\}}$ & $-1.33$ & $-1.42$ & $-1.64$ & $-4.84^{***}$ & $-3.31^{***}$ & $-2.86^{***}$ & $-0.98^{***}$ & $-1.02^{***}$ & $-1.43^{***}$ \\ 
   & $(1.01)$ & $(1.13)$ & $(1.10)$ & $(1.49)$ & $(1.14)$ & $(0.95)$ & $(0.14)$ & $(0.13)$ & $(0.13)$ \\ \vspace{5pt} 
  \hline \multicolumn{10}{c}{Panel B: Negative shocks} \\
 \hline
  & \multicolumn{3}{c}{BMI} & \multicolumn{3}{c}{FIT} & \multicolumn{3}{c}{OFI} \\ 
  \cmidrule(l){2-4} \cmidrule(l){5-7} \cmidrule(l){8-10} 
  & (1) & (2) & (3) & (4) & (5) & (6) & (7) & (8) & (9) \\ 
  $M_{\{|d_{n,t}| < \sigma\}}$ & $0.25$ & $-0.19$ & $-0.99$ & $2.47$ & $2.30^{*}$ & $2.74^{***}$ & $1.93^{***}$ & $1.75^{***}$ & $1.75^{***}$ \\ 
   & $(1.12)$ & $(1.27)$ & $(1.21)$ & $(1.62)$ & $(1.24)$ & $(1.02)$ & $(0.15)$ & $(0.11)$ & $(0.10)$ \\ 
  $M_{\{|d_{n,t}| \in [\sigma, 2\sigma]\}}$ & $0.98^{*}$ & $0.37$ & $0.15$ & $1.84^{*}$ & $1.86^{***}$ & $2.09^{***}$ & $1.44^{***}$ & $1.30^{***}$ & $1.26^{***}$ \\ 
   & $(0.59)$ & $(0.66)$ & $(0.59)$ & $(0.94)$ & $(0.69)$ & $(0.58)$ & $(0.12)$ & $(0.09)$ & $(0.08)$ \\ 
  $M_{\{|d_{n,t}| > 2 \sigma\}}$ & $0.29$ & $0.04$ & $-0.26$ & $2.08^{***}$ & $2.14^{***}$ & $2.27^{***}$ & $0.48^{***}$ & $0.43^{***}$ & $0.41^{***}$ \\ 
   & $(0.62)$ & $(0.61)$ & $(0.60)$ & $(0.69)$ & $(0.58)$ & $(0.54)$ & $(0.07)$ & $(0.06)$ & $(0.06)$ \\ \vspace{5pt} 
  Predictor controls & N & Y & Y & N & Y & Y & N & Y & Y \\ 
  Liquidity controls & N & N & Y & N & N & Y & N & N & Y \\ \vspace{5pt} 
  Obs & 4,681 & 4,681 & 4,681 & 234,705 & 234,705 & 234,705 & 306,871 & 306,871 & 306,871 \\ 
  $R^2$ & $0.074$ & $0.203$ & $0.248$ & $0.007$ & $0.073$ & $0.091$ & $0.026$ & $0.079$ & $0.102$ \\ 
  Marginal $R^2(d_{n,t})$ & $0.026$ & $0.026$ & $0.022$ & $0.007$ & $0.005$ & $0.004$ & $0.026$ & $0.017$ & $0.013$ \\ 
 \hline
  $M_{\{|d_{n,t}| \in [\sigma, 2\sigma]\}} - M_{\{|d_{n,t}| < \sigma\}}$ & $0.74$ & $0.56$ & $1.15$ & $-0.76$ & $-0.68$ & $-0.90$ & $-0.49^{***}$ & $-0.45^{***}$ & $-0.49^{***}$ \\ 
   & $(0.99)$ & $(1.01)$ & $(1.04)$ & $(0.99)$ & $(0.86)$ & $(0.78)$ & $(0.07)$ & $(0.07)$ & $(0.06)$ \\ 
  $M_{\{|d_{n,t}| > 2 \sigma\}} - M_{\{|d_{n,t}| \in [\sigma, 2\sigma]\}}$ & $-0.69$ & $-0.33$ & $-0.41$ & $0.39$ & $0.37$ & $0.27$ & $-0.96^{***}$ & $-0.87^{***}$ & $-0.86^{***}$ \\ 
   & $(0.81)$ & $(0.77)$ & $(0.73)$ & $(0.84)$ & $(0.74)$ & $(0.71)$ & $(0.09)$ & $(0.08)$ & $(0.07)$ \\ 
  $M_{\{|d_{n,t}| > 2 \sigma\}} - M_{\{|d_{n,t}| < \sigma\}}$ & $0.05$ & $0.23$ & $0.73$ & $-0.24$ & $-0.26$ & $-0.60$ & $-1.46^{***}$ & $-1.32^{***}$ & $-1.34^{***}$ \\ 
   & $(1.23)$ & $(1.32)$ & $(1.24)$ & $(1.42)$ & $(1.11)$ & $(0.97)$ & $(0.11)$ & $(0.09)$ & $(0.09)$ \\ \vspace{5pt} 
   \hline 
\end{tabular}


\endgroup
\end{adjustbox}
\end{table}

\paragraph{Controlled-specification highlights.}
In the most controlled specification, positive FIT multipliers fall from 6.43
in the smallest bin to 3.57 in the largest bin; positive OFI multipliers fall
from 2.99 to 1.56. Negative OFI multipliers fall from 1.75 to 0.41, with a
largest-minus-smallest difference of -1.34. Negative FIT falls from 2.74 to
2.27 between the smallest and largest bins, but the difference is imprecise.
BMI coefficients are not informative because the sign-specific BMI samples are
small and the estimates are noisy.

\end{document}
